\documentclass[twoside,final]{hcmut-report}

\usepackage{listings}
\usepackage{xcolor}
\usepackage{tcolorbox}

% ─── Python code listing style ────────────────────────────────────────────
\definecolor{codebg}{rgb}{0.96,0.96,0.96}
\definecolor{codekey}{rgb}{0.13,0.29,0.53}
\definecolor{codestr}{rgb}{0.58,0.00,0.10}
\definecolor{codecmt}{rgb}{0.40,0.40,0.40}
\lstdefinestyle{pystyle}{
    backgroundcolor=\color{codebg},
    basicstyle=\ttfamily\footnotesize,
    keywordstyle=\color{codekey}\bfseries,
    stringstyle=\color{codestr},
    commentstyle=\color{codecmt}\itshape,
    language=Python,
    frame=single,
    breaklines=true,
    showstringspaces=false,
    tabsize=2,
    captionpos=b,
    numbers=left,
    numberstyle=\tiny\color{codecmt}
}
\lstset{style=pystyle}

\coursename{Software Testing}
\reporttype{Project 3:}
\title{Data-driven Automation Testing and Non-functional Testing}
\advisor{& Bùi Hoài Thắng &}
\groupname{& 11 &}
\stuname{%
  & Lâm Minh Tùng Quân &  2352991 \\
  & Trương Xuân Sang &   2353045\\
  & Trương Gia Kỳ Nam &  2352787 \\
  & Nguyễn Hữu Phúc &  2352938 \\
  & Châu Anh Nhật  & 2352856\\
}

\begin{document}

\coverpage%

\tableofcontents
\pagebreak

% ============================================================
% CONTRIBUTION TABLE
% ============================================================
\begin{table}[ht]
    \centering
    \caption{Project Contribution and Workload Distribution}
    \vspace{0.2cm}
    \small
    \renewcommand{\arraystretch}{1.4}

    \begin{tabularx}{\textwidth}{|c|c|p{2.5cm}|>{\raggedright\arraybackslash}X|c|}
    \hline
    \textbf{No.} & \textbf{ID} & \textbf{Full Name} & \textbf{Description} & \textbf{\%} \\ \hline

    1 & 2352938 & Nguyễn Hữu Phúc &
    \begin{filelist}
        \item Feature 001: Admin Adds a New User -- Level 1 + Level 2 automation
        \item NFR-1: Performance Testing (Locust load testing)
        \item NFR-2: Security Testing (OWASP ZAP active scan)
    \end{filelist} & 20\% \\ \hline

    2 & 2352787 & Trương Gia Kỳ Nam &
    \begin{filelist}
       \item Feature 006: Teacher Sets Up a Quiz -- Level 1 + Level 2 automation
       \item NFR-1: Security Testing (OWASP ZAP active scan)
       \item NFR-2: Accessibility Testing (axe-core WCAG audit)
    \end{filelist} & 20\% \\ \hline

    3 & 2352856 & Châu Anh Nhật &
    \begin{filelist}
        \item Feature 005: User Creates Calendar Event -- Level 1 + Level 2 automation
        \item NFR-1: Performance Testing (Locust load testing)
        \item NFR-2: Accessibility Testing (axe-core WCAG audit)
    \end{filelist} & 20\% \\ \hline

    4 & 2353045 & Trương Xuân Sang &
    \begin{filelist}
       \item Feature 002: Admin Creates a New Course -- Level 1 + Level 2 automation
       \item Feature 003: Teacher Creates Assignment -- Level 1 + Level 2 automation
       \item NFR-1: Performance Testing (Locust load testing)
       \item NFR-2: Accessibility Testing (axe-core WCAG audit)
    \end{filelist} & 20\% \\ \hline

    5 & 2352991 & Lâm Minh Tùng Quân &
    \begin{filelist}
       \item Feature 004: Teacher Grades an Assignment -- Level 1 + Level 2 automation
       \item NFR-1: Security Testing (OWASP ZAP active scan)
       \item NFR-2: Accessibility Testing (axe-core WCAG audit)
    \end{filelist} & 20\% \\ \hline

    \end{tabularx}
\end{table}
\newpage

% ============================================================
% 1. GENERAL INFORMATION
% ============================================================
\section{General Information}

\subsection{Project Overview}
Project~\#3 extends the manual black-box test designs from Project~\#2 by transforming them into \textbf{data-driven automated test suites} written in Python with the Selenium WebDriver framework. In addition to the functional regression suites, this project introduces \textbf{non-functional testing} covering performance, security, and accessibility characteristics of the Moodle LMS demo environment.

The complete artefact set consists of:
\begin{itemize}
    \item \textbf{Level 1 automation} (data only from CSV) -- six suites, 169 parametric test methods.
    \item \textbf{Level 2 automation} (data \emph{and} locators / URLs from CSV) -- a unified suite, 169 parametric test methods.
    \item \textbf{Non-functional testing} -- seven Python files implementing five complementary NFR techniques.
    \item \textbf{Supporting tooling} -- a master PowerShell runner (\texttt{run\_all.ps1}), a Selenium-based test-data cleanup script (\texttt{cleanup\_moodle.py}), and a pinned \texttt{requirements.txt}.
\end{itemize}

\subsection{Target Application}
All automation runs against the dedicated MoodleCloud LMS instance hosted at
\begin{center}
    \textbf{\url{https://ihatetesting.moodlecloud.com/}}
\end{center}
which was established during Project~\#2 to bypass the instability of the public Moodle demo. The application instance, the administrator credentials, and the six feature flows under test remain identical to those of Project~\#2; only the testing technique has evolved.

\subsection{Test Environment}
To ensure reproducibility, every test session was conducted under the following pinned configuration:

\begin{table}[H]
\centering
\renewcommand{\arraystretch}{1.3}
\begin{tabularx}{\textwidth}{|l|X|}
\hline
\textbf{Component} & \textbf{Version / Detail} \\ \hline
Operating System & Windows 11 Pro 64-bit (build 22621) \\ \hline
Language & Python 3.10.11 (CPython) \\ \hline
Browser & Google Chrome 148.x (auto-managed ChromeDriver) \\ \hline
Test framework & Selenium 4.10+, pytest 9.0, unittest \\ \hline
Driver management & \texttt{webdriver-manager} 4.0 (auto ChromeDriver sync) \\ \hline
Data libraries & \texttt{pandas} 2.0, \texttt{openpyxl} 3.1 \\ \hline
Load testing & \texttt{locust} 2.20+ \\ \hline
Security scanning & OWASP ZAP 2.16 + \texttt{python-owasp-zap-v2.4} \\ \hline
Accessibility audit & \texttt{axe-selenium-python} 2.1 (axe-core engine) \\ \hline
Display resolution & 1920~$\times$~1080 (Full HD) at 100\% scale \\ \hline
\end{tabularx}
\caption{Pinned test environment}
\end{table}

\subsection{AI Tools Usage Declaration}
In strict compliance with academic-integrity guidelines, our team declares the following AI assistance:
\begin{itemize}
    \item \textbf{Google Gemini (Free Tier)} -- consulted for brainstorming boundary values, Python helper-function patterns, and \LaTeX{} formatting refinements.
    \item \textbf{Anthropic Claude (Sonnet/Opus)} -- consulted as a coding pair-programmer for debugging Selenium element-locator strategies, generating data-cleanup scripts, and refactoring the NFR test scaffolds.
    \item \textbf{Scope of use:} brainstorming, syntax assistance, log analysis, and refactoring suggestions only. The core test-design logic, automation strategy, and verification of every test result were manually performed and reviewed by the human team members.
\end{itemize}

% ============================================================
% 2. REPOSITORY LAYOUT
% ============================================================
\section{Repository Layout and Submission Package}

The submission archive is organised into three top-level folders matching the project rubric. Files for Level~1 and Level~2 are deliberately separated as required by the project description.

\begin{lstlisting}[language=bash, basicstyle=\ttfamily\footnotesize]
swe-testing/
|-- level1/                          # Level 1 - data only from CSV
|   |-- test_add_user_level1.py      # TC-001 (43 rows)
|   |-- test_data_tc001.csv
|   |-- test_course_level1.py        # TC-002 (27 rows)
|   |-- test_data_tc002.csv
|   |-- test_assign_level1.py        # TC-003 (27 rows)
|   |-- test_data_tc003.csv
|   |-- test_grade_level1.py         # TC-004 (17 rows)
|   |-- test_data_tc004.csv
|   |-- test_event_level1.py         # TC-005 (28 rows)
|   |-- test_data_tc005.csv
|   |-- test_quiz_level1.py          # TC-006 (27 rows)
|   `-- test_data_tc006.csv
|
|-- level2/                          # Level 2 - data + locators from CSV
|   |-- test_level2.py               # Single unified script (169 tests)
|   |-- test_data_tc001_level2.csv
|   |-- test_data_tc002_level2.csv
|   |-- test_data_tc003_level2.csv
|   |-- test_data_tc004_level2.csv
|   |-- test_data_tc005_level2.csv
|   `-- test_data_tc006_level2.csv
|
|-- non_functional/                  # Non-functional test suites
|   |-- test_nfr_01_login_perf_sec.py    # Locust + OWASP ZAP
|   |-- test_nfr_02_login_perf_a11y.py   # Locust + axe-core
|   |-- test_nfr_03_login_sec_a11y.py    # OWASP ZAP + axe-core
|   |-- test_nfr_04_quiz_perf_sec.py     # Locust + OWASP ZAP
|   |-- test_nfr_05_quiz_perf_a11y.py    # Locust + axe-core
|   `-- test_nfr_06_quiz_sec_a11y.py     # OWASP ZAP + axe-core
|   
|-- requirements.txt                 # Pinned dependencies
|-- run_all.ps1                      # Master PowerShell runner
|-- cleanup_moodle.py                # Test-data cleanup utility
|-- README.md                        # Full usage guide
`-- report.pdf                       # This report
\end{lstlisting}

% ============================================================
% 3. LEVEL 1 -- DATA-DRIVEN AUTOMATION
% ============================================================
\section{Level 1 -- Data-driven Automation Testing}

\subsection{Design Principles}
Level~1 implements the data-driven testing pattern as defined in the project description: each test case is parametrised by a row of input data and an \texttt{expected\_result} value read from a CSV file, while the locators (element identifiers) remain hard-coded inside the Python script.

The architecture for every Level~1 test file follows the same template:
\begin{enumerate}
    \item A module-level CSV reader (\texttt{csv.DictReader}) loads all rows from \texttt{test\_data\_tcXXX.csv}.
    \item A factory function (\texttt{\_make\_test(row)}) returns a closure that becomes the body of one \texttt{unittest} test method.
    \item A \texttt{setattr} loop attaches one such method per CSV row at import time, so \texttt{pytest} collects them as independent tests.
    \item Each test method calls the same \texttt{\_fill\_and\_submit(row)} helper, then asserts that the observed outcome matches \texttt{row["expected\_result"]}.
\end{enumerate}

\subsection{Per-feature Implementation Details}

This section documents each of the six features automated in Level~1. Test
case purposes are grouped by the black-box technique that generated them
(BVA, ECP, DT, UC) -- mirroring the categorisation used in the Project~\#2
report. Each CSV row in the data file corresponds to one parametric test
method named \texttt{test\_TC\_XXX\_YYY} at pytest collection time.

% ─────────────────────────────────────────────────────────────────────────
\subsubsection{Feature 001 -- Admin Adds a New User (43 test cases)}

\textbf{Target URL:} \url{https://ihatetesting.moodlecloud.com/user/editadvanced.php?id=-1}\\
\textbf{Pre-condition:} Logged in as administrator on the Moodle site.\\
\textbf{Data file:} \texttt{level1/test\_data\_tc001.csv} (43 rows).\\
\textbf{CSV schema:} \texttt{test\_case\_id, username, password, firstname, lastname, email, expected\_result}.

\textbf{Special handling:}
\begin{itemize}
    \item Moodle marks \texttt{\#id\_newpassword} as \texttt{readonly}; the script bypasses this through JavaScript injection (\texttt{removeAttribute}~+~\texttt{dispatchEvent}).
    \item The sentinel value \texttt{\_\_generate\_\_} in the password column ticks the \emph{Generate password and notify user} checkbox instead of typing a password.
\end{itemize}
\textbf{Outcome detection:} success $\Leftrightarrow$ ``Changes saved'' substring in page source; fail $\Leftrightarrow$ any \texttt{[id\^{}='id\_error\_']} element rendered.

\textbf{Test-case purpose grouped by technique:}
\begin{table}[H]
\centering
\renewcommand{\arraystretch}{1.2}
\begin{tabularx}{\textwidth}{|c|l|X|c|}
\hline
\textbf{TC range} & \textbf{Method} & \textbf{Boundary / scenario tested} & \textbf{Exp.} \\ \hline
TC-001-001 & BVA & Username (nom 17 chars) -- valid user, all fields nominal & success \\ \hline
TC-001-002 & BVA & Username (min$-$) empty string -- missing required field & fail \\ \hline
TC-001-003..004 & BVA & Username (min 1 char, min$+$ 2 chars) -- boundary lower & success \\ \hline
TC-001-005 & BVA & Password (min$-$ 7 chars) -- below 8-char minimum & fail \\ \hline
TC-001-006..010 & BVA & Password length boundaries (8, 9, 127, 128, 129 chars) & success/fail \\ \hline
TC-001-011 & BVA & First name (min$-$) empty -- required field missing & fail \\ \hline
TC-001-012..015 & BVA & First name boundaries (1, 2, 99, 100 chars) & success \\ \hline
TC-001-016 & BVA & Last name (min$-$) empty -- required field missing & fail \\ \hline
TC-001-017..020 & BVA & Last name boundaries (1, 2, 99, 100 chars) & success \\ \hline
TC-001-021..024 & ECP & Invalid username formats: uppercase, space, special char, ``admin'' (duplicate) & fail \\ \hline
TC-001-025..028 & DT & Password complexity rules: missing digit / upper / lower / non-alphanumeric & fail \\ \hline
TC-001-029..031 & ECP & Invalid email formats: no @, missing domain, embedded space & fail \\ \hline
TC-001-032 & UC & Happy-path: all fields nominal -- baseline success scenario & success \\ \hline
TC-001-033 & UC & \texttt{\_\_generate\_\_} sentinel -- alternate flow with auto-password & success \\ \hline
TC-001-034 & UC & Duplicate username ``admin'' -- second create attempt after TC-001-024 & fail \\ \hline
TC-001-035..043 & ECP/DT & Additional password \& email edge combinations & varies \\ \hline
\end{tabularx}
\caption{TC-001 test case purpose by group (43 rows total)}
\end{table}

% ─────────────────────────────────────────────────────────────────────────
\subsubsection{Feature 002 -- Admin Creates a New Course (27 test cases)}

\textbf{Target URL:} \url{https://ihatetesting.moodlecloud.com/course/edit.php?category=1}\\
\textbf{Data file:} \texttt{level1/test\_data\_tc002.csv} (27 rows).\\
\textbf{CSV schema:} \texttt{test\_case\_id, fullname, shortname, end\_date\_enabled, end\_date\_offset\_days, end\_date\_offset\_years, numsections, expected\_result}.

\textbf{Special handling:} End-date day / month / year drop-downs are set through JavaScript helpers (\texttt{sS()} for select boxes, \texttt{ens()}~/~\texttt{dis()} for the \texttt{id\_enddate\_enabled} checkbox).\\
\textbf{Outcome detection:} success $\Leftrightarrow$ ``Announcements'' substring (course landing page); fail $\Leftrightarrow$ any \texttt{[id\^{}='id\_error\_']} validation element.

\textbf{Test-case purpose grouped by technique:}
\begin{table}[H]
\centering
\renewcommand{\arraystretch}{1.2}
\begin{tabularx}{\textwidth}{|c|l|X|c|}
\hline
\textbf{TC range} & \textbf{Method} & \textbf{Boundary / scenario tested} & \textbf{Exp.} \\ \hline
TC-002-001 & BVA & Full name (nom 20 chars) + shortname \texttt{sn001\_xxxxxxxx} -- nominal success & success \\ \hline
TC-002-002 & BVA & Full name (min$-$) empty -- required field missing & fail \\ \hline
TC-002-003..004 & BVA & Full name boundaries (1, 2 chars) & success \\ \hline
TC-002-005..006 & BVA & Full name boundaries (253, 254 chars max) & success \\ \hline
TC-002-007 & BVA & Short name (min$-$) empty -- required field missing & fail \\ \hline
TC-002-008..009 & BVA & Short name boundaries (1, 2 chars) & success \\ \hline
TC-002-010..011 & BVA & Short name boundaries (254, 255 chars max) & success \\ \hline
TC-002-012..013 & BVA & End date chronology (-1 day, today vs start) -- invalid order & fail \\ \hline
TC-002-014..017 & BVA & End date boundaries (+1 day .. +12 years from start) & success \\ \hline
TC-002-018 & ECP & Duplicate short name -- second create attempt with same shortname & fail \\ \hline
TC-002-019..023 & DT & Course format \texttt{weekly}/\texttt{topics} + end-date enabled/disabled combinations & success \\ \hline
TC-002-024..025 & UC & Happy paths (all valid + end-date disabled variants) & success \\ \hline
TC-002-026 & UC & Duplicate short name as intentional second attempt -- mirrors TC-002-018 & fail \\ \hline
TC-002-027 & UC & Past end date -- date-conflict exception scenario & fail \\ \hline
\end{tabularx}
\caption{TC-002 test case purpose by group (27 rows total)}
\end{table}

% ─────────────────────────────────────────────────────────────────────────
\subsubsection{Feature 003 -- Teacher Creates Assignment (27 test cases)}

\textbf{Target URL:} \url{https://ihatetesting.moodlecloud.com/course/modedit.php?add=assign&type&course=425&sectionid=2116&return=0&beforemod=0}\\
\textbf{Role switch:} \texttt{setUpClass()} navigates to \texttt{/course/switchrole.php?id=1\&switchrole=-1} to assume the Teacher role before any assignment tests run.\\
\textbf{Data file:} \texttt{level1/test\_data\_tc003.csv} (27 rows).\\
\textbf{CSV schema:} \texttt{test\_case\_id, name, gradepass, duedate\_enabled, duedate\_offset\_days, duedate\_offset\_years, cutoff\_offset\_days, cutoff\_offset\_years, submission\_file, submission\_onlinetext, expected\_result}.

\textbf{Special handling:} Three date fields (\texttt{allowsubmissionsfromdate}, \texttt{duedate}, \texttt{cutoffdate}) are populated via a single \texttt{execute\_script()} JS helper that also enables / disables their checkboxes per CSV row. Submission-type checkboxes (\texttt{file}, \texttt{onlinetext}) are only changed when the CSV column value is not \texttt{default}.\\
\textbf{Outcome detection:} ``Announcements'' substring after redirect $\Rightarrow$ success; any \texttt{[id\^{}='id\_error\_']} $\Rightarrow$ fail.

\textbf{Test-case purpose grouped by technique:}
\begin{table}[H]
\centering
\renewcommand{\arraystretch}{1.2}
\begin{tabularx}{\textwidth}{|c|l|X|c|}
\hline
\textbf{TC range} & \textbf{Method} & \textbf{Boundary / scenario tested} & \textbf{Exp.} \\ \hline
TC-003-001 & BVA & Assignment name (nom 29 chars) -- baseline valid create & success \\ \hline
TC-003-002 & BVA & Name (min$-$) empty -- required field missing & fail \\ \hline
TC-003-003..004 & BVA & Name boundaries (1, 2 chars lower bound) & success \\ \hline
TC-003-005..010 & BVA & Grade-to-pass boundaries (0, 0.01, 9.99, 10, 10.01, negative) & success/fail \\ \hline
TC-003-011..016 & BVA & Due date relative to allow-submission date (today, past, future, +10/11/12y) & success/fail \\ \hline
TC-003-017..022 & BVA & Cut-off date relative to due date (-1 day, due, due+1, due+10y, ..) & success/fail \\ \hline
TC-003-023..025 & DT & Submission-type combinations: file only, online only, both, neither & success \\ \hline
TC-003-026..027 & UC & Happy path + empty-name exception scenario & success/fail \\ \hline
\end{tabularx}
\caption{TC-003 test case purpose by group (27 rows total)}
\end{table}

% ─────────────────────────────────────────────────────────────────────────
\subsubsection{Feature 004 -- Teacher Grades an Assignment (17 test cases)}

\textbf{Target URL:} \url{https://ihatetesting.moodlecloud.com/mod/assign/view.php?id=1195&action=grader&userid=2}\\
\textbf{Pre-condition:} Logged in as administrator, switched to Teacher role.\\
\textbf{Data file:} \texttt{level1/test\_data\_tc004.csv} (17 rows).\\
\textbf{CSV schema:} \texttt{test\_case\_id, grade, expected\_result}.

\textbf{Special handling:} The grade input is a React-controlled component, so a regular \texttt{send\_keys} does not register. The script uses the native \texttt{HTMLInputElement.prototype.value} setter and dispatches \texttt{input}~/~\texttt{change}~/~\texttt{blur} events to make React re-render. After submit, a sentinel DOM element \texttt{\#\_\_test\_marker[data-has-error]} is injected via JavaScript so the outcome can be read after the page navigates away.

\textbf{Test-case purpose grouped by technique:}
\begin{table}[H]
\centering
\renewcommand{\arraystretch}{1.2}
\begin{tabularx}{\textwidth}{|c|l|X|c|}
\hline
\textbf{TC range} & \textbf{Method} & \textbf{Boundary / scenario tested} & \textbf{Exp.} \\ \hline
TC-004-001..004 & BVA & Grade boundaries lower-bound (0, 0.01, 99.99, 100) & success \\ \hline
TC-004-005..006 & BVA & Grade out of range ($-1$, $100.01$) -- below min / above max & fail \\ \hline
TC-004-007..010 & ECP & Non-numeric grades: letters, special chars, empty (feedback only) & fail/success \\ \hline
TC-004-011..017 & BVA/UC & Decimal precision edge values + happy/notify combinations & varies \\ \hline
\end{tabularx}
\caption{TC-004 test case purpose by group (17 rows total)}
\end{table}

% ─────────────────────────────────────────────────────────────────────────
\subsubsection{Feature 005 -- User Creates Calendar Event (28 test cases)}

\textbf{Target URL:} \url{https://ihatetesting.moodlecloud.com/calendar/view.php?view=month}\\
\textbf{Data file:} \texttt{level1/test\_data\_tc005.csv} (28 rows).\\
\textbf{CSV schema:} \texttt{test\_case\_id, name, duration\_type, minutes, until\_offset\_days, repeat, expected\_result}.

\textbf{Special handling:} Three duration modes (\texttt{none}, \texttt{minutes}, \texttt{until}) map to radio buttons (values 0, 1, 2). \texttt{minutes} additionally fills a numeric field; \texttt{until} sets a future date through the \texttt{sS()} helper offset by \texttt{until\_offset\_days}. The modal submit button is located via \texttt{//div[@role='dialog']//button[@data-action='save']}.\\
\textbf{Outcome detection:} modal dialog closes within 8~s + ``Calendar'' substring in page source.

\textbf{Test-case purpose grouped by technique:}
\begin{table}[H]
\centering
\renewcommand{\arraystretch}{1.2}
\begin{tabularx}{\textwidth}{|c|l|X|c|}
\hline
\textbf{TC range} & \textbf{Method} & \textbf{Boundary / scenario tested} & \textbf{Exp.} \\ \hline
TC-005-001..004 & BVA & Event title boundaries (nom, empty, 1, 2 chars) & success/fail \\ \hline
TC-005-005..010 & BVA & Duration in minutes boundaries (0, 1, 2, 9999998, 9999999, 10000000) & success/fail \\ \hline
TC-005-011..015 & BVA & ``Until'' date offset ($-1$, 0, $+1$, $+364$, $+365$ days) & success/fail \\ \hline
TC-005-016..018 & ECP & Invalid minutes inputs (1.5 decimal, alphabetic, negative) & fail \\ \hline
TC-005-019..021 & DT & Duration mode combinations (none, minutes only, until only) & success \\ \hline
TC-005-022..024 & DT & Repeat checkbox enabled in each duration mode & success \\ \hline
TC-005-025..028 & UC & Happy path + exception flows (empty title, negative duration) & varies \\ \hline
\end{tabularx}
\caption{TC-005 test case purpose by group (28 rows total)}
\end{table}

% ─────────────────────────────────────────────────────────────────────────
\subsubsection{Feature 006 -- Teacher Sets Up a Quiz (27 test cases)}

\textbf{Target URL:} \url{https://ihatetesting.moodlecloud.com/course/modedit.php?add=quiz&type&course=426&sectionid=2121&return=0&beforemod=0}\\
\textbf{Role switch:} \texttt{setUpClass()} switches to Teacher role on course 426, with editing mode enabled.\\
\textbf{Data file:} \texttt{level1/test\_data\_tc006.csv} (27 rows).\\
\textbf{CSV schema:} \texttt{test\_case\_id, name, timeclose\_enabled, close\_offset\_days, close\_offset\_years, timelimit\_enabled, timelimit\_number, gradepass, expected\_result}.

\textbf{Special handling:} The \texttt{expected\_result} column distinguishes four failure sub-modes (\texttt{fail\_grade}, \texttt{fail\_time}, \texttt{fail\_date}, \texttt{fail\_name}) identified by the specific \texttt{id\_error\_*} element rendered after submit. This gives more precise assertions than a generic \texttt{fail}.

\textbf{Test-case purpose grouped by technique:}
\begin{table}[H]
\centering
\renewcommand{\arraystretch}{1.2}
\begin{tabularx}{\textwidth}{|c|l|X|c|}
\hline
\textbf{TC range} & \textbf{Method} & \textbf{Boundary / scenario tested} & \textbf{Exp.} \\ \hline
TC-006-001..004 & BVA & Grade-to-pass boundaries (5, 9.99, 10, 10.01) -- above max grade & success/fail\_grade \\ \hline
TC-006-005..010 & BVA & Time-limit boundaries ($-1$, 0, 1, 998, 999, 1000) & success/fail\_time \\ \hline
TC-006-011..016 & BVA & Close-date boundaries ($-1$ day, today, $+1$ day, $+10$/$11$/$12$ years) & success/fail\_date \\ \hline
TC-006-017 & ECP & Empty quiz name -- required field missing & fail\_name \\ \hline
TC-006-018..024 & DT & Timing combinations (open + close + limit enabled/disabled rules R1..R5) & success \\ \hline
TC-006-025..026 & UC & Composite exception flows (past close date $+$ grade $>$ max) & fail \\ \hline
TC-006-027 & UC & Empty name in happy-path context -- final exception scenario & fail\_name \\ \hline
\end{tabularx}
\caption{TC-006 test case purpose by group (27 rows total)}
\end{table}

\subsection{Sample Code Skeleton}
\begin{lstlisting}[caption={Generic Level~1 test-factory pattern},label={lst:lvl1-factory}]
import csv, os, unittest
from selenium import webdriver
from selenium.webdriver.common.by import By

CSV_PATH = os.path.join(os.path.dirname(__file__), "test_data_tc001.csv")

def _load_rows():
    with open(CSV_PATH, newline="", encoding="utf-8") as f:
        return list(csv.DictReader(f))

class TestAddUserLevel1(unittest.TestCase):
    @classmethod
    def setUpClass(cls):
        cls.driver = webdriver.Chrome()
        cls._login()

    def _fill_and_submit(self, row):
        # Locators are hardcoded here (Level 1 contract)
        self.driver.find_element(By.ID, "username").send_keys(row["username"])
        # ... fill other fields ...
        self.driver.find_element(By.ID, "id_submitbutton").click()

    def _get_outcome(self):
        if "Changes saved" in self.driver.page_source:
            return "success"
        if self.driver.find_elements(By.CSS_SELECTOR, "[id^='id_error_']"):
            return "fail"
        return "unknown"

def _make_test(row):
    def test_method(self):
        self._fill_and_submit(row)
        self.assertEqual(self._get_outcome(),
                         row["expected_result"].strip().lower())
    test_method.__name__ = f"test_{row['test_case_id'].replace('-', '_')}"
    return test_method

for _r in _load_rows():
    setattr(TestAddUserLevel1, f"test_{_r['test_case_id'].replace('-','_')}",
            _make_test(_r))
\end{lstlisting}

% ============================================================
% 4. LEVEL 2 -- LOCATORS + URLS FROM CSV
% ============================================================
\section{Level 2 -- Locators and URLs Externalised}

\subsection{Design Principles}
Level~2 fully realises the project description's strongest requirement: \emph{``testing data and testing items such as site URLs and elements like text fields and buttons are provided to the script.''} In other words, the Python source contains \textbf{zero site-specific values}: every URL, every locator strategy, and every locator value is read from CSV columns.

\subsection{Locator Resolution Helper}
A single helper \texttt{loc(row, prefix)} translates a pair of CSV columns into a Selenium locator tuple:

\begin{lstlisting}[caption={Locator resolution from CSV},label={lst:loc-helper}]
from selenium.webdriver.common.by import By

BY_MAP = {
    "id": By.ID, "name": By.NAME, "css": By.CSS_SELECTOR,
    "xpath": By.XPATH, "link text": By.LINK_TEXT,
    "partial link text": By.PARTIAL_LINK_TEXT,
    "tag name": By.TAG_NAME, "class name": By.CLASS_NAME,
}

def loc(row, prefix):
    """Resolve (By.X, "value") from columns <prefix>_locator_type and
       <prefix>_locator_value."""
    return (BY_MAP[row[f"{prefix}_locator_type"].strip().lower()],
            row[f"{prefix}_locator_value"].strip())
\end{lstlisting}

Every interactive element thus contributes two CSV columns, e.g.\ for the \texttt{username} input on TC-001 Level~2:
\begin{center}
    \texttt{username\_locator\_type = id, \quad username\_locator\_value = id\_username}
\end{center}

\subsection{Unified Class Architecture}
All six TC suites live in a single file \texttt{level2/test\_level2.py}. A shared base class \texttt{\_BaseLevel2} owns the browser lifecycle, login, cookie-banner dismissal, and a session-recovery hook (\texttt{setUp()} probes \texttt{driver.current\_url} and re-creates the driver if it has died):

\begin{center}
\renewcommand{\arraystretch}{1.3}
\begin{tabular}{|l|l|c|}
\hline
\textbf{Suite class} & \textbf{Feature} & \textbf{Test count} \\ \hline
\texttt{TestCreateUserLevel2} & TC-001 -- Add User & 43 \\ \hline
\texttt{TestCreateCourseLevel2} & TC-002 -- Create Course & 27 \\ \hline
\texttt{TestAssignLevel2} & TC-003 -- Create Assignment & 27 \\ \hline
\texttt{TestGradeLevel2} & TC-004 -- Grade Assignment & 17 \\ \hline
\texttt{TestCalendarEventLevel2} & TC-005 -- Calendar Event & 28 \\ \hline
\texttt{TestQuizSetupLevel2} & TC-006 -- Quiz Setup & 27 \\ \hline
\multicolumn{2}{|r|}{\textbf{Total}} & \textbf{169} \\ \hline
\end{tabular}
\end{center}

\subsection{CSV Schema Example (Excerpt -- TC-001 Level~2)}
\begin{lstlisting}[language=, basicstyle=\ttfamily\scriptsize, frame=single, breaklines=true]
test_case_id,site_url,login_url_suffix,new_user_url,username,password,
firstname,lastname,email,expected_result,
username_locator_type,username_locator_value,
firstname_locator_type,firstname_locator_value,
lastname_locator_type,lastname_locator_value,
email_locator_type,email_locator_value,
save_btn_locator_type,save_btn_locator_value

TC-001-001,https://ihatetesting.moodlecloud.com,/login/index.php,
https://ihatetesting.moodlecloud.com/user/editadvanced.php?id=-1,
usr001_uuuuuuuuuu,XuanSang12@,ffffffffff,llllllllll,usr001_nom@test.com,
success,id,id_username,id,id_firstname,id,id_lastname,id,id_email,
id,id_submitbutton
\end{lstlisting}

\subsection{Test Execution Summary -- Level 2}
A complete Level~2 run on the pinned environment yields:
\begin{center}
\fbox{\parbox{0.8\textwidth}{\centering
\texttt{====== 169 passed in 2323.73s (0:38:43) ======}
}}
\end{center}
i.e.\ all 169 parametric tests pass green; total wall-clock time $\approx$~39~minutes.

% ============================================================
% 5. NON-FUNCTIONAL TESTING
% ============================================================
\section{Non-functional Testing}

The project description requires \emph{at least two non-functional requirements per student}, with each NFR described in terms of \textbf{testing type, testing approach, and testing tool}. The team identified \textbf{three complementary NFR categories} (matching the Python-native recommendations in \texttt{NFR.md}):
\begin{itemize}
    \item \textbf{Performance} -- load testing with \textbf{Locust} (NFR-B below).
    \item \textbf{Security} -- DAST scanning with \textbf{OWASP ZAP} Python API (NFR-D below).
    \item \textbf{Accessibility} -- WCAG audit with \textbf{axe-selenium-python} (NFR-E below).
\end{itemize}

For 5 team members $\times$ 2 NFRs each $=$ 10 NFR slots, the three categories are paired so that every student demonstrates two distinct NFR types on their owned feature (see Table~\ref{tab:per-student-nfr} below). Each of the 6 NFR files implements two of the three techniques applied to either the Login or the Quiz Setup flow, providing pairwise coverage across the team.

\subsection{NFR Catalogue}

\subsubsection{NFR-1 -- Performance: Load Testing}
\textbf{Type:} Performance / concurrent load (multi-user, throughput, p95 latency, failure rate).\\
\textbf{Approach:} Define a virtual user as a Python class deriving from \texttt{HttpUser}; each \texttt{@task} method represents one user interaction (visit login page, submit form, open quiz form, etc.). Locust spawns dozens of these virtual users in parallel and aggregates timing statistics in real time through its web UI on port~8089.\\
\textbf{Tool:} Locust 2.43 (\url{https://locust.io}).\\
\textbf{Files:} \texttt{test\_nfr\_01\_login\_perf\_sec.py}, \texttt{test\_nfr\_02\_login\_perf\_a11y.py}, \texttt{test\_nfr\_04\_quiz\_perf\_sec.py}, \texttt{test\_nfr\_05\_quiz\_perf\_a11y.py}.

\subsubsection{NFR-2 -- Security: Active Scanning with OWASP ZAP}
\textbf{Type:} Security / Dynamic Application Security Testing (DAST).\\
\textbf{Approach:} OWASP ZAP runs as a daemon on \texttt{127.0.0.1:8080}; the test scripts drive it through the official Python API (\texttt{python-owasp-zap-v2.4}). Each test performs the canonical ZAP workflow: (1) Spider the target URL to enumerate attack surface; (2) Active Scan injects 100+ canned payloads (SQLi, XSS, path-traversal, command injection, etc.) and observes responses; (3) Aggregate findings into \emph{High / Medium / Low / Informational} risk levels. The test fails if any High-risk alert is reported.\\
\textbf{Tool:} OWASP ZAP 2.16 + \texttt{python-owasp-zap-v2.4} 0.0.21.\\
\textbf{Files:} \texttt{test\_nfr\_01\_login\_perf\_sec.py}, \texttt{test\_nfr\_03\_login\_sec\_a11y.py}, \texttt{test\_nfr\_04\_quiz\_perf\_sec.py}, \texttt{test\_nfr\_06\_quiz\_sec\_a11y.py}.

\subsubsection{NFR-3 -- Accessibility: WCAG Audit with axe-core}
\textbf{Type:} Accessibility / WCAG 2.1 conformance.\\
\textbf{Approach:} While a Selenium session has a target page open, \texttt{axe-core} (the industry-standard accessibility engine) is injected into the page through \texttt{axe-selenium-python}. The engine audits the live DOM against $\sim$90~rules covering colour-contrast, ARIA labelling, keyboard focus order, heading hierarchy, and so on. The test fails if any violation of impact \emph{``critical''} or \emph{``serious''} is reported. A full JSON report is persisted alongside the test for evidence.\\
\textbf{Tool:} \texttt{axe-selenium-python} 2.1 (wraps \texttt{axe-core} 4.x).\\
\textbf{Files:} \texttt{test\_nfr\_02\_login\_perf\_a11y.py}, \texttt{test\_nfr\_03\_login\_sec\_a11y.py}, \texttt{test\_nfr\_05\_quiz\_perf\_a11y.py}, \texttt{test\_nfr\_06\_quiz\_sec\_a11y.py}.

\subsection{NFR Coverage Matrix}
Each of the six NFR files combines exactly two of the three NFR techniques (Locust load, ZAP security, axe accessibility) applied to one of the two main user flows (Login or Quiz Setup). The pairwise coverage is therefore demonstrated twice in independent contexts.

\begin{table}[H]
\centering
\renewcommand{\arraystretch}{1.3}
\begin{tabularx}{\textwidth}{|l|l|c|c|c|}
\hline
\textbf{File} & \textbf{Flow} & \textbf{Performance} & \textbf{Security} & \textbf{Accessibility} \\ \hline
\texttt{test\_nfr\_01\_login\_perf\_sec.py}   & Login & \checkmark (Locust) & \checkmark (ZAP) & -- \\ \hline
\texttt{test\_nfr\_02\_login\_perf\_a11y.py}  & Login & \checkmark (Locust) & -- & \checkmark (axe) \\ \hline
\texttt{test\_nfr\_03\_login\_sec\_a11y.py}   & Login & -- & \checkmark (ZAP) & \checkmark (axe) \\ \hline
\texttt{test\_nfr\_04\_quiz\_perf\_sec.py}    & Quiz  & \checkmark (Locust) & \checkmark (ZAP) & -- \\ \hline
\texttt{test\_nfr\_05\_quiz\_perf\_a11y.py}   & Quiz  & \checkmark (Locust) & -- & \checkmark (axe) \\ \hline
\texttt{test\_nfr\_06\_quiz\_sec\_a11y.py}    & Quiz  & -- & \checkmark (ZAP) & \checkmark (axe) \\ \hline
\end{tabularx}
\caption{Pairwise NFR coverage matrix}
\end{table}

\subsection{Per-Student NFR Assignment}
\begin{table}[H]
\centering
\renewcommand{\arraystretch}{1.3}
\begin{tabularx}{\textwidth}{|p{3.6cm}|X|X|}
\hline
\textbf{Student} & \textbf{NFR-1} & \textbf{NFR-2} \\ \hline
Nguyễn Hữu Phúc      & Performance -- Locust load testing  & Security -- OWASP ZAP active scan \\ \hline
Trương Xuân Sang     & Performance -- Locust load testing  & Accessibility -- axe-core WCAG audit \\ \hline
Lâm Minh Tùng Quân   & Security -- OWASP ZAP active scan   & Accessibility -- axe-core WCAG audit \\ \hline
Châu Anh Nhật        & Performance -- Locust load testing  & Accessibility -- axe-core WCAG audit \\ \hline
Trương Gia Kỳ Nam    & Security -- OWASP ZAP active scan   & Accessibility -- axe-core WCAG audit \\ \hline
\end{tabularx}
\caption{Each student owns two distinct NFR types (project requirement)}
\label{tab:per-student-nfr}
\end{table}

% ============================================================
% 6. SUPPORTING TOOLING
% ============================================================
\section{Supporting Tooling}

\subsection{Master Test Runner -- \texttt{run\_all.ps1}}
A single PowerShell entry point orchestrates the entire test pipeline (install $\to$ ZAP daemon lifecycle $\to$ sequenced suite execution). It exposes ten non-interactive modes plus an interactive menu:

\begin{lstlisting}[language=PowerShell, basicstyle=\ttfamily\footnotesize]
.\run_all.ps1                 # Interactive menu
.\run_all.ps1 -Mode setup     # Install dependencies, probe Chrome/ZAP
.\run_all.ps1 -Mode smoke     # 1-test sanity check (~30 s)
.\run_all.ps1 -Mode level1    # All 169 Level 1 tests (~30 min)
.\run_all.ps1 -Mode level2    # All 169 Level 2 tests (~30 min)
.\run_all.ps1 -Mode nfr-old   # Baseline NFR (Performance + Security)
.\run_all.ps1 -Mode nfr-new   # 6 NFR files (Locust + ZAP + axe)
.\run_all.ps1 -Mode locust    # Pick a Locust load-test file
.\run_all.ps1 -Mode cleanup   # Delete Moodle test data
.\run_all.ps1 -Mode all       # Everything end-to-end (~1 h 40 m)
\end{lstlisting}

The script auto-starts and tears down the ZAP daemon (even on Ctrl-C), so the operator never has to babysit a hidden process.

\subsection{Test-data Cleanup -- \texttt{cleanup\_moodle.py}}
Because the data-driven tests intentionally do \emph{not} clean up after themselves (sequential dependencies between rows -- e.g.\ \texttt{TC-001-034} validates that creating a duplicate \texttt{admin} fails, so the previous test \texttt{TC-001-024} must have created it), accumulated test rows would eventually pollute the Moodle instance. The cleanup utility addresses this between full runs:

\begin{itemize}
    \item \textbf{Users} -- bulk delete on \texttt{/admin/user/user\_bulk.php} matching pattern \texttt{usr\textbackslash{}d+\_*}, mirroring the Katalon TC-001-Cleanup flow.
    \item \textbf{Courses} -- search-page enumeration matching pattern \texttt{fn*}~/~\texttt{sn*}, deleted through the two-step \texttt{/course/delete.php} confirmation flow.
    \item \textbf{Activities} -- per-course enumeration (course~425 for TC-003 assignments, course~426 for TC-006 quizzes) with names matching test patterns, deleted through a JavaScript \texttt{fetch}~+~\texttt{DOMParser} fast-path that submits the Moodle confirm form directly.
    \item \textbf{Calendar events} -- enumeration from \texttt{/calendar/view.php?view=upcoming}, deleted through \texttt{/calendar/delete.php}.
\end{itemize}

Protected entities (built-in \texttt{admin} user, the team's administrator account, courses 425 and 426 themselves) are excluded by a hard-coded allowlist.

\subsection{Reproducibility Manifest -- \texttt{requirements.txt}}
\begin{lstlisting}[language=, basicstyle=\ttfamily\footnotesize, frame=single]
selenium>=4.10           # WebDriver core
webdriver-manager>=4.0   # Auto-managed ChromeDriver
pandas>=2.0              # CSV / Excel loading
openpyxl>=3.1            # XLSX support
pytest>=7.4              # Test runner
locust>=2.20             # Performance load testing
python-owasp-zap-v2.4>=0.0.21  # Security scanning API
axe-selenium-python>=2.1.6     # Accessibility auditing
\end{lstlisting}
A single \texttt{pip install -r requirements.txt} restores the entire stack on a clean machine.

% ============================================================
% 7. TEST EXECUTION RESULTS
% ============================================================
\section{Test Execution Results}

\subsection{Aggregate Pass / Fail Summary}
The final clean run on the pinned environment produces the following results:

\begin{table}[H]
\centering
\renewcommand{\arraystretch}{1.3}
\begin{tabularx}{\textwidth}{|l|c|c|c|X|}
\hline
\textbf{Suite} & \textbf{Total} & \textbf{Pass} & \textbf{Fail} & \textbf{Wall-clock} \\ \hline
Level 1 -- TC-001 & 43 & 43 & 0 & $\sim$8 min \\ \hline
Level 1 -- TC-002 & 27 & 27 & 0 & $\sim$5 min \\ \hline
Level 1 -- TC-003 & 27 & 27 & 0 & $\sim$5 min \\ \hline
Level 1 -- TC-004 & 17 & 17 & 0 & $\sim$3 min \\ \hline
Level 1 -- TC-005 & 28 & 28 & 0 & $\sim$5 min \\ \hline
Level 1 -- TC-006 & 27 & 27 & 0 & $\sim$5 min \\ \hline
\textbf{Level 1 total} & \textbf{169} & \textbf{169} & \textbf{0} & $\sim$31 min \\ \hline
\textbf{Level 2 total} & \textbf{169} & \textbf{169} & \textbf{0} & 38 min 43 s \\ \hline
NFR files (Locust + ZAP + axe) & 16 & 16 & 0 & 39 min 6 s \\ \hline
\end{tabularx}
\caption{Final clean-run summary}
\end{table}

\subsection{Notable Findings During Development}

\paragraph{TC-001-024 -- order-dependent test.}
This test attempts to create a user with username \texttt{admin}, expecting a duplicate-username failure. On a fresh MoodleCloud instance the built-in \texttt{admin} did not exist, so the first run unexpectedly succeeded. After running once, the user named \texttt{admin} is now permanently present, restoring the expected \texttt{fail} outcome on every subsequent run.

\paragraph{Duplicate course short names.}
After multiple iterations of the test suite, course shortnames \texttt{sn001\_ssssssss}, \texttt{sn003\_ssssssss}, \texttt{sn004\_ssssssss} accumulated in the database, causing duplicate-key failures on row 001/003/004 of TC-002. The CSV was updated to suffix \texttt{\_xxxxxxxx} (and the intentional-duplicate row TC-002-026 was kept in sync), restoring green status.

\paragraph{Course / section ID migration.}
The hard-coded course IDs from Project~\#2 (\texttt{course=141}, \texttt{sectionid=695} for TC-003; \texttt{course=152}, \texttt{sectionid=750} for TC-006) referenced resources that were later cleaned up on the MoodleCloud admin side. New dedicated courses were created (\texttt{425~=~TC\_003\_MAIN}, \texttt{426~=~TC\_004\_MAIN}), the active section~IDs were discovered automatically (\texttt{2116} and \texttt{2121} respectively), and all five affected files (three Python sources plus two Level~2 CSVs plus four NFR files) were updated atomically.

\paragraph{Encoding pitfall -- UTF-8 BOM.}
The Windows PowerShell~5.1 cmdlet \texttt{Set-Content -Encoding utf8} silently prepends a Byte-Order Mark (BOM) to the file. CSV files re-saved this way caused \texttt{csv.DictReader} to read the first column header as \texttt{\textbackslash{}ufefftest\_case\_id}, breaking every test-factory call. The team's fix uses \texttt{[System.IO.File]::WriteAllText} with an explicit \texttt{UTF8Encoding(\$false)} encoder to guarantee BOM-free output.

% ============================================================
% 8. CONCLUSION
% ============================================================
\section{Conclusion}
Project~\#3 successfully transformed the manual black-box test designs of Project~\#2 into a robust, reproducible, two-level data-driven automation suite covering all six target features (169 test methods per level, 338 in total), plus seven non-functional test files covering five complementary NFR techniques (response-time SLA, load testing, input-validation security, DAST scanning, and WCAG accessibility).

Throughout this project we gained practical experience in:
\begin{itemize}
    \item \textbf{Test architecture} -- balancing the trade-off between locator hard-coding (Level 1, simpler) and full CSV externalisation (Level 2, more flexible but more verbose CSV schemas).
    \item \textbf{Cross-browser reliability} -- using \texttt{webdriver-manager} to insulate the suite from ChromeDriver / Chrome version drift, and JavaScript-based DOM manipulation to bypass overlay interception, React-controlled inputs, and read-only attributes.
    \item \textbf{Non-functional methodology} -- integrating three distinct industry-standard tools (Locust, OWASP ZAP, axe-core) into a unified Python testing harness, and documenting each according to the project rubric's \emph{type / approach / tool} triad.
    \item \textbf{Test-data lifecycle} -- recognising that sequential dependencies between data rows are intentional, and that cleanup must therefore run \emph{between} suite runs rather than \emph{after every test}.
    \item \textbf{Operational tooling} -- packaging a complex multi-stage test pipeline behind a single \texttt{run\_all.ps1} entry point, complete with daemon-process lifecycle management and graceful fallback when optional tools (e.g.\ ZAP) are not installed.
\end{itemize}

The final clean run achieves \textbf{169/169 green on Level~1}, \textbf{169/169 green on Level~2}, and \textbf{all non-functional assertions passing} on the pinned environment, demonstrating that the Moodle LMS demo instance meets the documented functional, performance, security, and accessibility expectations of this academic test scope.

\end{document}
